\section{Prescribing the full automorphism group}\label{co:symmetry}
A graph automorphism is an adjacency-preserving permutation whose inverse
also preserves adjacency. A finite group is given by a multiplication
table on $[m]$ with $m\ge1$, an identity $1_G$ and an inverse table.
The group equations are finite tests. For each $a$ the permutation
$L_a(x)=ax$ has inverse $L_{a^{-1}}$ and satisfies $L_aL_b=L_{ab}$.
Evaluation at the identity distinguishes the translations. This already
proves Cayley's permutation representation theorem.

For a subset $S$ closed under inversion and excluding the identity,
join $x$ to $xs$ for each $s\in S$. The resulting simple undirected
Cayley graph has degree $|S|$ by cancellation. Left translations give
a transitive group of automorphisms. Its component at $1_G$ is exactly
the subgroup generated by $S$: a path spells a product of generators,
and any product gives a walk by its successive partial products. More
explicitly, the reachable mask at $1_G$ is closed under multiplication
and inversion by concatenating generator words and reversing them with
inverted letters. Translating it gives all the other components as left
cosets. Reachability needs paths of length at most $m-1$, so the component
description also computes subgroup membership.

\begin{theorem}[Finite Frucht realization]\leavevmode\label{co:frucht}
One can construct a connected simple undirected graph on at most
$(m+2m^2)(2m+2)$ vertices whose full automorphism group is isomorphic
to the given group. Both directions of the isomorphism are explicit.
\end{theorem}
\begin{proof}
Introduce a vertex $v_x$ for each group element $x$. For each ordered
pair $x,g$ use distinct new vertices $a_{x,g},b_{x,g}$ and the edges
of the walk
\[
 v_x\;--\; a_{x,g}\;--\; b_{x,g}\;--\; v_{xg}.
\]
When $g=1_G$ these three edges form a triangle. All are ordinary
edges of a simple graph, including when $m=1$. Call the vertices
introduced so far the core. At each core vertex attach a fresh path
by one endpoint. Its length is one at $v_x$, $2+2g$ at $a_{x,g}$,
and $3+2g$ at $b_{x,g}$, using $g\in[m]$ as its numerical label.
There are $m+2m^2$ core vertices and each attached path adds at most
$2m+1$ vertices. All core vertices can reach a $v_x$, and the walks
from $v_{1_G}$ reach every $v_x$. The graph is connected.

In the finished graph, core vertices are exactly those of degree at
least three: at each $a$ or $b$ there are two distinct core neighbours
and the attached path, and each $v_x$ has $2m$ distinct core neighbours.
All remaining vertices lie on paths and have degree one or two.
Each core vertex has exactly one noncore neighbour, and following the
unique attached path recovers its length. It follows that any automorphism
preserves the three types, including $g$ on both inner types. Write
$v_x\mapsto v_{\pi(x)}$. There is a unique neighbour of type $a_{*,g}$
at each $v_x$, so $a_{x,g}\mapsto a_{\pi(x),g}$. Its unique neighbour
of type $b_{*,g}$ then forces $b_{x,g}\mapsto b_{\pi(x),g}$. The
remaining core adjacency gives $\pi(xg)=\pi(x)g$. Substituting
$x=1_G$ proves $\pi=L_{\pi(1_G)}$. Images of all core vertices have now
been determined, and distances along their rooted paths force the images
of all other vertices.

Conversely each $L_z$ acts on the three core types by replacing $x$ by
$zx$ and preserving $g$. Associativity preserves the terminal incidence
at $v_{xg}$, and sending each path position to the same position on the
image path extends this to an automorphism with inverse indexed by $z^{-1}$.
The resulting automorphisms compose by group multiplication. Reading
the image of $v_{1_G}$ gives the inverse function on the full
automorphism group. Given the graph order bound $M$, all permutation and
inverse tests scan words of length at most $M$ with entries below $M$.
Core types and path lengths are recovered in at most $M$ steps.
\end{proof}
